\begin{table}[htb!]
\centering
\caption{Priority counties that socio-economic screening would miss. A priority county is one ranked in both the highest tercile of expected tropical cyclone (TC) damage and the lowest tercile of construction capacity across the analysis sample ($n = 685$ counties); these are places where high storm risk coincides with a thin local rebuilding sector. A county is considered flagged by a socio-economic index when it falls in that index's most concerning tercile, that is, the highest tercile of the FEMA National Risk Index (NRI) Social Vulnerability score (SVI) or the lowest tercile of the NRI Community Resilience score (Resil.). The 13 counties listed are priority counties in neither of those terciles, so a screen built on either NRI index would overlook them; they are the counties where the building-permit capacity screen adds information beyond socio-economic indicators. Counties are sorted by number of housing units. Pop.\ \% of state is the percentage of the state population living in the county. Permits per month is the construction-capacity proxy (average monthly building permits), and permits per 1000 units normalises it by housing stock (sample median 0.046). NRI SVI and Community Resilience are county percentile scores from 0 to 100; a higher SVI indicates greater social vulnerability and a higher Resilience score indicates greater resilience.}
\label{tab:divergent_counties}
\begin{tabular}{llrrrrrr}
\hline
County & State & Housing & Pop.\ \% & Permits & Permits per & NRI & NRI \\
 & & units & of state & per month & 1000 units & SVI & Resil. \\
\hline
Cameron & Louisiana & 3,582 & 0.12 & 2.25 & 0.628 & 45 & 99 \\
Union & Florida & 5,293 & 0.07 & 3.75 & 0.708 & 39 & 32 \\
Seminole & Georgia & 5,969 & 0.09 & 0.33 & 0.056 & 48 & 33 \\
Grant & Louisiana & 9,228 & 0.48 & 0.42 & 0.045 & 58 & 39 \\
Allen & Louisiana & 9,880 & 0.49 & 3.83 & 0.388 & 61 & 39 \\
Appling & Georgia & 10,125 & 0.17 & 0.33 & 0.033 & 74 & 49 \\
Brantley & Georgia & 10,489 & 0.17 & 4.08 & 0.389 & 34 & 31 \\
Pike & Alabama & 18,021 & 0.66 & 3.50 & 0.194 & 74 & 58 \\
Ware & Georgia & 19,934 & 0.34 & 3.75 & 0.188 & 78 & 81 \\
Lamar & Mississippi & 27,228 & 2.17 & 3.75 & 0.138 & 56 & 56 \\
St. Mary & Louisiana & 29,762 & 1.06 & 3.08 & 0.104 & 76 & 88 \\
Jones & Mississippi & 33,761 & 2.27 & 1.92 & 0.057 & 78 & 41 \\
Victoria & Texas & 43,846 & 0.31 & 3.58 & 0.082 & 59 & 64 \\
\hline
\end{tabular}
\end{table}
